\documentclass[10pt]{article}
\usepackage[a4paper,margin=25mm]{geometry}
\usepackage{amsmath,amssymb,amsthm,mathrsfs,longtable,hyperref}
\emergencystretch=3em
\newtheorem{theorem}{Theorem}[section]
\newtheorem{lemma}[theorem]{Lemma}
\newtheorem{proposition}[theorem]{Proposition}
\newtheorem{corollary}[theorem]{Corollary}
\newcommand{\PSH}{\operatorname{PSH}}
\newcommand{\FS}{\operatorname{FS}}
\newcommand{\MA}{\operatorname{MA}}
\newcommand{\vol}{\operatorname{vol}}
\newcommand{\OO}{\mathcal O}
\newcommand{\E}{\mathcal E}
\newcommand{\an}{\mathrm{an}}
\newcommand{\dd}{dd^c}
\title{Absolute model estimates and a conditional attainment mechanism\\
\large A standalone incomplete paper for the unchanged original P16.5}
\author{Research draft from the two specified mathematical grants}
\date{6 October 2026}
\begin{document}
\maketitle
\begin{abstract}
We prove a raw finite-model absolute-value contraction, retaining all central
multiplicities, and compare the midpoint energy defect with the finite Hodge
distance. We independently reconstruct the uniform finite-model Taylor
estimate from full section lattices, positive deficient measures, and algebraic
volume differentiation. A smallness condition keeps both signed perturbations
inside the big cone. We then give a conditional original-carrier continuation
of the proposed arbitrary-target approximation: an absolute midpoint estimate,
transferred with the missing-rank error retained, rules out energy jumps at
model near-maximizers. Finite divisorial supporting planes suffice for actual
attainment, and divisorial probes give strict separation. This continuation
requires the explicitly listed genuine scalar, compactness, and moving-output
foundations; those foundations are not proved in their entirety here. Full
unbounded mixed independence and strict Borel comparison are also not proved.
Thus this paper does not solve original P16.5 and has canonical result
\texttt{candidate=false}. No zero mass is probability-normalized.
\end{abstract}

\section{Exact scope and mathematical authority}
Throughout the sought positive-volume domain is every algebraically closed,
trivially valued characteristic-zero field $k$, every smooth integral
projective $k$-variety $X$ of dimension $n\geq1$, and every actual big rational
line bundle $L\in\operatorname{Pic}(X)\otimes\mathbb Q$.
Write $K=X^\an$, $\theta=c_1(L)$, $V=\vol_X(L)>0$, and let $v_0$ denote
the generic trivial valuation. No countability, finite-generation, toric,
skeletal, or target-support restriction is added. All supported semivaluations
remain on the original compact Hausdorff carrier $K$. Canonical section norms
are at most one and are one at $v_0$ for each nonzero section.
An actual bundle identity is required for section multiplication; numerical
classes enter only ampleness and intersections. All rational tensor powers
and divisor coefficients are cleared before sections or ideals are used.

The only caller grants are the entire supplied G0, in its \emph{actual ample}
rational-bundle domain and exact retained field and geometric hypotheses,
and G1, in its exact resolved smooth equivariant relative-generation signature
and its \emph{one} generically finite FS net. We do not regard the other supplied
bodies or ledgers as premises. In particular the uniform variation calculation
below is a new reconstruction, whereas the scalar/moving package in
Section~\ref{package} is a list of unresolved antecedents. All uses of its
assertions are explicitly conditional.

Ordinary algebraic background consists of intersection theory on smooth
projective models, characteristic-zero equivariant principalization and SNC
resolution, nef positivity, Kodaira decomposition, high-power ample generation,
normal proper birational section descent, section-volume asymptotics, and the
ordinary algebraic positive-product and volume results of BFJ. These are not
nonample NA equation grants. Complete pertinent BFJ and finite BJ proof units
were retrieved independently; the source appendix gives literal coordinates.
We do not assert a recursive new proof of resolution, Fujita approximation,
or algebraic Hodge index.

The cone convention retained here is the literal single-net convention.
An FS scalar on an actual $B$ has the form
\[
 h=m^{-1}\max_j\{\log\|s_j\|_{\rm triv}+\lambda_j\},\qquad
 0\ne s_j\in H^0(X,mB),\quad\lambda_j\in\mathbb R.
\]
Common zeros and supported $-\infty$ poles are allowed. $\PSH(\theta)$
consists of all-point decreasing limits of \emph{one} such net on varying
actual bundles $B_i$ with $c_1(B_i)\to\theta$ and finite generic endpoint.
This definition alone does not give decreasing closure of a twice-closed
family. The intended reference is $R=P_\theta(0)$ and the intended energy is
the outer full-section energy specified in Section~\ref{package}. We never put
a polar physical section function into $\E^1$ merely by factoring its divisor.

One elementary embedding uses G0 without nonample closure: if $\Omega-L$ is
ample, a tail of a defining net has $\Omega-B_i$ ample. Tensor its sections
with a weight-zero generating system of a cleared power of that actual
difference. The scalar does not change, since a generator is a unit at the
center of every semivaluation. To make the stages finite, append generators
of $\Omega$ at weights $-mN$ and use the product net $\max(h_i,-N)$.
This is a genuine single decreasing ample net. G0 gives finite divisorial
values and all-point order determination for its endpoint. Thus for each
fixed rational divisorial $z$ there is a finite $C_z$ with
\begin{equation}\label{coordinate-bound}
 -C_z\leq u(z)\leq0
 \quad\text{for every normalized }u\in\PSH(\theta).
\end{equation}
Here normalization is $u(v_0)=\sup_Ku=0$. G0 compactness gives the bound
in the one actual ample cone. No equality of ample and big energies is used.
Rational divisorial scalings suffice for order determination: along a
divisorial ray $t\operatorname{ord}_E$, every defining FS stage is a finite
maximum of affine functions of $t>0$. Its decreasing endpoint is convex
and finite by the ample embedding, hence continuous on $(0,\infty)$.
An inequality on rational scalings therefore extends to all real scalings,
after which G0's divisorial order determination applies. This uses neither
countability of the field nor equality only on a skeletal carrier.

\section{Absolute-value contraction with every multiplicity}
Let $Y\to Z\times\mathbb P^1$ be a smooth SNC equivariant model, product
away from zero, with $\dim Z=n$. Write its full central fiber as
$F=\sum_i b_iE_i$, $b_i>0$. Let $N_1,\ldots,N_{n-1}$ be actual rational
relatively nef model bundles. For a vertical rational Cartier divisor
$D=\sum_i d_iE_i$, put $a_i=d_i/b_i$, and let $f_D$ be its scalar model
function. When $n=1$ the product of the $N_j$ is empty.
Denote $\Gamma=N_1\cdots N_{n-1}$.

\begin{lemma}[Graph identity]\label{graph}
With $c_{ij}=b_ib_j(E_iE_j\Gamma)$ for $i<j$, one has
\begin{equation}\label{graph-eq}
 -D^2\Gamma=\sum_{i<j}c_{ij}(a_i-a_j)^2,\qquad c_{ij}\geq0.
\end{equation}
For every rational continuous model function $f$,
\begin{equation}\label{abs-contract}
 \|\,|f|\,\|_\Gamma^2\leq\|f\|_\Gamma^2,
 \qquad\|f_D\|_\Gamma^2=-D^2\Gamma.
\end{equation}
The same contraction holds for a rational Lipschitz piecewise linear
function $P(f)$ of Lipschitz constant one, provided it is represented by
the corresponding resolved model function.
\end{lemma}
\begin{proof}
The line bundle $\OO(F)$ restricts trivially to each $E_i$. Hence
$\sum_j b_j(E_iE_j\Gamma)=0$ for every $i$. For $i\ne j$, $E_iE_j$ is
an effective SNC stratum, and the relatively nef $N_l$ restrict to nef
bundles on this projective vertical stratum. Thus its intersection is
nonnegative. Expanding $D^2\Gamma$, replacing its diagonal terms with
the preceding zero-row identities, and collecting unordered pairs proves
\eqref{graph-eq}. This derivation includes the diagonal coefficients,
the $b_i$, and every exceptional component.

Resolve the absolute value \emph{before} applying the graph identity.
Choose integral $q$ clearing $D$ and rational $C$ sufficiently large that
$q(CF+D)$ and $q(CF-D)$ are effective integral vertical Cartier divisors.
Principalize the invariant ideal
\[
 \OO(-q(CF+D))+\OO(-q(CF-D)).
\]
On a dominating SNC model it is $\OO(-G)$, and $CF-G/q$ represents
$|f_D|$: its semivaluative coefficient is
$C-\min(C+f_D,C-f_D)=|f_D|$.
The ideal is the unit ideal off the central fiber; no horizontal poles
are introduced. On every central component of this refined model the
normalized coefficient of the new divisor is precisely $|a_i|$ for the
pulled $D$. Pull back all $N_l$ and $F$, retaining their new central
multiplicities. Apply \eqref{graph-eq} on this model and use
$||a_i|-|a_j||\leq|a_i-a_j|$ term by term. Projection formula preserves
the old right quadratic form. This proves \eqref{abs-contract}.
The last assertion has exactly the same proof using
$|P(a_i)-P(a_j)|\leq|a_i-a_j|$ after resolving its rational maxima and
minima. No abstract chain rule or unspecified Dirichlet form is assumed.
\end{proof}

\section{Midpoint defect and a raw absolute moment}
On one actual nef rational generic bundle $A$ on a smooth projective
$n$-fold, let $\xi,\eta,\psi$ be bounded rational relatively nef model
scalars. Put $T_\xi=(A,\xi)$, $T_0=(A,0)$, $V_A=A^n$, and
\[
 e_A(\xi)=T_\xi^{n+1}/(n+1),\quad
 J_A(\xi)=\xi(v_0)V_A-e_A(\xi),\quad
 d_A(\xi,\eta)=\max_{0\leq j<n}
   \|\xi-\eta\|_{T_\xi^jT_\eta^{n-1-j}}^2.
\]
These are finite raw intersection quantities. Relatively nef model top
measures are positive with mass $V_A$ and coefficient
$b_E(T_\xi|_E)^n$ at $v_E=b_E^{-1}\operatorname{ord}_E|_{k(Z)}$.
They follow from G0 by a fixed-model ample perturbation and finite
polynomial passage, or from the finite intersection calculus directly.
No unbounded nef or big equation is invoked.

Two telescoping sums, with $S=(0,\xi-\eta)$, give
\begin{equation}\label{correct}
 T_\eta^{n+1}-T_\xi^{n+1}+(n+1)ST_\eta^n
 =-\sum_{j=0}^{n-1}(n-j)S^2T_\xi^jT_\eta^{n-1-j}.
\end{equation}
Indeed the term with index $j$ occurs in the inner telescoping sum for
$k=j+1,\ldots,n$. Consequently
\begin{equation}\label{J-formula}
 J_A(\xi)=\frac1{n+1}\sum_{j=0}^{n-1}(n-j)
       \|\xi\|_{T_\xi^jT_0^{n-1-j}}^2\asymp_n d_A(\xi,0).
\end{equation}
The orientation $j+1$ with these retained factors is not used.

\begin{lemma}[Midpoint defect]\label{midpoint}
For $\delta_A(\xi,\eta)=e_A((\xi+\eta)/2)
 -(e_A(\xi)+e_A(\eta))/2$, one has
\begin{equation}\label{midbound}
 \frac{n}{4^{n+1}}d_A(\xi,\eta)
 \leq\delta_A(\xi,\eta)\leq\frac n8d_A(\xi,\eta).
\end{equation}
\end{lemma}
\begin{proof}
Let $p(t)=e_A(\eta+t(\xi-\eta))$. Multilinearity gives
\[
 -p''(t)=n\sum_{j=0}^{n-1}\binom{n-1}{j}
          t^j(1-t)^{n-1-j}Q_j,\quad
 Q_j=\|\xi-\eta\|_{T_\xi^jT_\eta^{n-1-j}}^2\geq0.
\]
The one-dimensional Green formula for linear interpolation gives
$\delta_A=\int_0^1 k(t)(-p''(t))dt$ with
$k(t)=\min(t,1-t)/2$. Its integral is $1/8$, giving the upper bound.
On $[1/4,3/4]$, $k\geq1/8$ and each
$t^j(1-t)^{n-1-j}\geq4^{1-n}$. Every coefficient of $Q_j$ is therefore
at least $n/4^{n+1}$. This gives the lower bound for their maximum.
For $n=1$ the exact identity is $\delta_A=Q_0/8$.
\end{proof}

\begin{lemma}[Finite slot changes]\label{slots}
Writing $\beta=2^{1-n}$, one has for relatively nef model potentials
$\xi,\eta,\rho_1,\ldots,\rho_{n-1}$
\begin{equation}\label{slotbound}
 \|\xi-\eta\|_{T_{\rho_1}\cdots T_{\rho_{n-1}}}^2
 \leq C_n d_A(\xi,\eta)^\beta
       \max\{J_A(\xi),J_A(\eta),J_A(\rho_l)\}^{1-\beta}.
\end{equation}
\end{lemma}
\begin{proof}
We record the finite proof to specify its scope. Expanding the average
$\tau=(\xi+\eta)/2$ gives
$d_A(\xi,\eta)\asymp_n\|\xi-\eta\|_{T_\tau^{n-1}}^2$.
The tangent remainder in \eqref{correct}, divided by $n+1$, is
comparable to $d_A$ and is convex in its first potential by concavity
of the finite energy. Thus $d_A$ is quasi-convex in that potential.

Fix another potential $\rho$. Set $a=d_A(\xi,\eta)$,
$b=\max(d_A(\xi,\rho),d_A(\eta,\rho))$, $f=\xi-\eta$, and
$b_j=\|f\|_{T_\rho^jT_\tau^{n-1-j}}^2$.
The average formula gives $b_0\leq C_na$; the seminorm triangle
inequality gives $b_{n-1}\leq4b$.
Replacing one slot and writing $(0,f)=T_\xi-T_\eta$ gives
\[
 b_{j+1}-b_j=-(0,f)^2(0,\rho-\tau)
                       T_\rho^jT_\tau^{n-2-j}.
\]
Cauchy--Schwarz in the two nef products obtained from $T_\xi,T_\eta$,
and $T_\xi,T_\eta\leq2T_\tau$ in these quadratic forms, bound this
by $C_n\sqrt{b_jd_A(\tau,\rho)}\leq C'_n\sqrt{b_jb}$.
The quadratic forms are nonnegative by Lemma~\ref{graph}; hence their
polarizations satisfy Cauchy--Schwarz. If $a\leq b$, induction gives
$b_j\leq C_na^{2^{-j}}b^{1-2^{-j}}$. If $a>b$, use $b_{n-1}\leq4b$.
This proves the single repeated-slot estimate.

For a third potential $\zeta$, apply that estimate to both differences
$\xi-\zeta$ and $\zeta-\eta$ in the average $T_\tau^{n-1}$.
Quasi-convexity and the average formula yield
$d\leq C_nb^\beta\max(d,b)^{1-\beta}$, where
$d=d_A(\xi,\eta)$ and $b=\max(d_A(\xi,\zeta),d_A(\eta,\zeta))$.
If $d\geq b$ divide by $d^{1-\beta}$; otherwise the conclusion is
immediate. Thus $d_A$ satisfies a dimension-only quasi-triangle bound.
Average all the desired $\rho_l$. Expanding its $(n-1)$st power
dominates each nonnegative mixed quadratic form, and the repeated-slot
bound, quasi-convexity, quasi-triangle bound, and \eqref{J-formula}
give \eqref{slotbound}. When $n=1$, there is one empty-slot seminorm
and the assertion is immediate with $\beta=1$.
\end{proof}

\begin{theorem}[Absolute finite moment]\label{absolute-finite}
Set $\alpha=2^{-n}$ and
$B_* =\max(J_A(\xi),J_A(\eta),J_A(\psi))$.
Then
\begin{equation}\label{absolute-finite-eq}
 \int|\xi-\eta|\,d\MA_{A,\rm raw}(\psi)
 \leq |\xi(v_0)-\eta(v_0)|V_A
      +C_n\delta_A(\xi,\eta)^\alpha
         J_A(\psi)^{1/2}B_*^{1/2-\alpha}.
\end{equation}
This assertion includes $V_A=0$, with no division by volume.
\end{theorem}
\begin{proof}
Let $g=|\xi-\eta|$ and resolve it as in Lemma~\ref{graph}. Since the
trivial metric has measure $V_A\delta_{v_0}$, telescoping from $j=0$ gives
\[
 \int g\,d(\MA_A(\psi)-\MA_A(0))
  =\sum_{j=0}^{n-1}(0,g)(0,\psi)T_\psi^jT_0^{n-1-j}.
\]
Bound each term by Cauchy--Schwarz. The first seminorm is at most the
one for $\xi-\eta$ by \eqref{abs-contract}, and \eqref{slotbound}
controls it. The second is at most $C_nJ_A(\psi)^{1/2}$ by
\eqref{J-formula}. Use \eqref{midbound} and sum the $n$ terms.
This proves \eqref{absolute-finite-eq}. Constants have zero quadratic
form, so no normalization changes the estimate. If $V_A=0$, positivity
and mass force the top measure to be zero already. Zero-factor cases
are read by continuity; for $n=1$ the last factor has exponent zero
and is omitted.
\end{proof}

\section{Independent uniform finite-model variation}\label{variation-section}
In this section $L$ is any of the retained actual big rational bundles.
Let $Y\to X\times\mathbb P^1$ be a smooth equivariant modification,
product away from zero, $F=\sum b_EE$, and $D$ an actual vertical
rational Cartier divisor. Write $f=f_D$ and $d=n+1$.
The full relative section lattice, for $R_m=H^0(X,mL)$, is
\begin{equation}\label{lattice}
 \Lambda_m=\bigoplus_j t^{-w_{mj}}k[t]s_{mj},\qquad
 ma\leq w_{mj}\leq mb
\end{equation}
when $aF\leq D\leq bF$. Indeed divisor comparison gives
$t^{-ma}R_m[t]\subset\Lambda_m\subset t^{-mb}R_m[t]$;
equivariance grades this lattice and an adapted basis of its finite
homogeneous flag gives \eqref{lattice}. Proper normal birational section
descent retains the \emph{whole} original $R_m$. No finite generation
of the section ring is assumed. Gluing to the product lattice at infinity gives
\begin{equation}\label{dimformula}
 h^0(Y,m(\rho_X^*L+D+cF))
          =\sum_j\max(w_{mj}+mc+1,0).
\end{equation}
For sufficiently large rational $c$, this is $(mc+1)\dim R_m+\sum_jw_{mj}$.
Section-volume limits therefore give an intrinsic eventual intercept
\begin{equation}\label{intercept}
 \vol_Y(\rho_X^*L+D+cF)/d=cV+\mathscr E(f).
\end{equation}
The full-section lattice shows that $\mathscr E$ is invariant under
model refinement, monotone, $V$-translation equivariant, and
$V$-Lipschitz for the original sup norm. Ordinary volume-root concavity,
applied to \eqref{intercept} and expanded at $c\to\infty$, makes it
concave on rational model functions. It has a unique Lipschitz extension
to $C(K)$ by actual ample model density from G0.

BFJ volume differentiation on this \emph{big algebraic} $(n+1)$-fold gives,
for large $c$,
\begin{equation}\label{finite-mu}
 \mu_D=\sum_E b_E\langle(\rho_X^*L+D+cF)^n\rangle\cdot E\,\delta_{v_E},
 \quad \mu_D(K)=V,\quad D\mathscr E_f(h)=\int h\,d\mu_D.
\end{equation}
Every coefficient is nonnegative by the volume derivative in the effective
$E$ direction. Each is independent of sufficiently large $c$, by
differentiating \eqref{intercept} in $D+sE$ on a common eventual half-line.
The derivative in $F$ gives the mass. Refinement is compatible: an
effective exceptional prime $G$ has $p_*\OO(rG)=\OO$ for $r\geq0$;
adding it to a pulled bundle does not change sections. The derivative
coefficient there is zero. Projection and pullback retain the other
coefficients. Equivalently the original functional and its model
derivatives are unchanged. This constructs finite model measures, not
an unbounded original operator.

\begin{lemma}[Positive deficient measures]\label{deficient}
Suppose $f_D\leq0$ and $\mathscr E(f_D)\geq-B$, $B\geq0$.
For every $c>B/V$, the class $\alpha_D(c)=\rho_X^*L+D+cF$ is big, and
\begin{equation}\label{deficient-eq}
 0\leq\mu_{D,c}:=\sum_Eb_E\langle\alpha_D(c)^n\rangle\cdot E\,\delta_{v_E}
 \leq\mu_D,\qquad (\mu_D-\mu_{D,c})(K)\leq B/c.
\end{equation}
\end{lemma}
\begin{proof}
Set $g_D(c)=\vol(\alpha_D(c))/d$. Formula \eqref{dimformula} gives
a convex function of $c$, with $g_D(0)=0$, slope bounded by $V$,
eventual expression $cV+\mathscr E(f_D)$, and
$g_D(c)\geq cV+\mathscr E(f_D)$. For $c>B/V$ this is positive,
so the actual bundle is big. At that point BFJ differentiation gives
$g_D'(c)=\mu_{D,c}(K)$. Convexity gives
$g_D'(c)\geq g_D(c)/c\geq V-B/c$, and its eventual slope gives
$g_D'(c)\leq V$.

The coefficient assertion requires more than this mass bound. The full
filtered-space trace in \eqref{dimformula} gives
\[
 T_D(c):=\mathscr E(f_D)+cV-g_D(c)
       =\lim_m\frac{n!}{m^{n+1}}\sum_j\min(w_{mj}+mc,0).
\]
Terms arising from the $+1$ in \eqref{dimformula} tend to zero, since
$\dim R_m=O(m^n)$. Increasing an effective vertical divisor increases
each ordered weight of its full lattice: this follows from inclusion
of every filtration space, or its dimension distribution. Thus each
finite trace of the increasing function $\min(x+mc,0)$ increases.
The same is true of $T_D(c)$. Its derivative in each effective $E$
is nonnegative and equals the coefficient of $\mu_D-\mu_{D,c}$
divided by $b_E$. This proves coefficientwise positivity.
Together with the mass estimate it proves \eqref{deficient-eq}.
\end{proof}

\begin{theorem}[Uniform finite-model Taylor estimate]\label{taylor}
Fix a rational original model direction $h$, a cost bound $B$, and
$c>B/V$. There exist $K_c<\infty$ and $t_c>0$, independent of the
nonpositive obstacle $f_D$ and of its refinement, such that
\begin{equation}\label{taylor-eq}
 |\mathscr E(f_D+th)-\mathscr E(f_D)-t\int h\,d\mu_D|
 \leq K_ct^2/2+|t|\|h\|_\infty(2B+2V\|h\|_\infty)/c
\end{equation}
whenever $\mathscr E(f_D)\geq-B$ and $|t|<t_c$.
\end{theorem}
\begin{proof}
Represent $h$ by a signed vertical Cartier $H$ on one fixed smooth
determination $Z$. Choose an actual ample $T$ there, and refine to carry
each $D$, always pulling $T,H$ back. Let
$\gamma=\rho_X^*L+cF$, $\alpha=\alpha_D(c)$,
$\kappa=d(cV-B)>0$, and $M=d\langle\gamma^n\rangle\cdot T$.
We have $\alpha\leq\gamma$ in psef order and $\vol(\alpha)\geq\kappa$.
While $\alpha-rT$ stays big, volume differentiation and positive-product
monotonicity bound its volume loss by $rM$. Thus
$\alpha-r_0T$ is big for $r_0=\kappa/[2(M+1)]$: a first earlier exit
would contradict volume continuity and the remaining positive volume.
The constants are unchanged under refinement, since their tests live on $Z$.

Choose $a>0$ with $aT\pm H$ nef and put $\lambda=a/r_0$.
Then $\lambda\alpha\pm H$ are psef. For $|s|\lambda<1/2$, positivity
and monotonicity sandwich
\[
 (1-\lambda|s|)^n\langle\alpha^n\rangle
 \leq\langle(\alpha+sH)^n\rangle
 \leq(1+\lambda|s|)^n\langle\alpha^n\rangle.
\]
To pair with $H$, decompose it into the \emph{two fixed nef tests}
$H_1=aT+H$, $H_2=aT$. Their pairings are bounded by
$\langle\gamma^n\rangle\cdot H_l$, fixed before the model varies.
The sandwich implies a common Lipschitz bound on the derivative of
$\vol(\alpha+sH)/d$. Integrating gives
\[
 |g_{D+tH}(c)-g_D(c)-t\int h\,d\mu_{D,c}|\leq K_ct^2/2.
\]
No signed test is inserted into an effective-order inequality.

Put $H_0=\|h\|_\infty$. At parameter $s$, use the nonpositive obstacle
$f_D+sh-|s|H_0$ and time parameter $c_s=c+|s|H_0$. Its cost is at
most $B_s=B+2V|s|H_0$, by the actual Lipschitz and constant laws.
Crucially choose $t_c\leq1$ and, if $H_0>0$,
$t_c<(c-B/V)/(2H_0)$. Then $c_s>B_s/V$ for both signs throughout
the integration path. The rough fixed bound $B+2VH_0$ alone would
not justify this big-cone condition when $c$ is near $B/V$.
Its deficient class is exactly $\alpha+sH$. Its full derivative measure
is unchanged by the scalar shift. Lemma~\ref{deficient} bounds the
positive derivative tail by $B_s/c_s\leq(B+2VH_0)/c$.
Integrate the $h$-derivative along this path and replace the initial
deficient derivative by $\mu_D$, at extra cost $|t|H_0B/c$.
The resulting sum is \eqref{taylor-eq}. The finite-model energy is $C^1$
on the path, since one sufficiently large time twist realizes its
intercept in a common neighborhood of actual big classes. This proves
the integration for both signs. When $H_0=0$ there is nothing to prove.
\end{proof}

Theorem~\ref{taylor} verifies the \emph{finite-model} variation ingredient.
It does not itself construct the scalar outer energy on every genuine base,
prove its envelope identity, construct a weakly continuous unbounded output,
or identify weighted moving products. Those distinctions are retained next.

\section{Explicit unresolved original-carrier antecedents}\label{package}
The rest of the positive-volume conclusions are conditional on the following
package S--M. It is stated for the \emph{entire} retained domain of Section 1,
with the literal single-net cone; no narrower space is substituted.
We do not claim to have proved the entire package independently.
\begin{enumerate}
\item[S1.] That cone is convex, finite-max closed and decreasing closed at
finite generic endpoints. Its normalized space is compact in the topology
of pointwise convergence at rational divisorial valuations. Each upper
bounded usc obstacle with a genuine generically finite minorant has a
genuine all-point envelope. The minimal reference is $R=P_\theta(0)$,
$R(v_0)=0$, and every normalized genuine $u$ satisfies $u\leq R\leq0$.
\item[S2.] The scalar energy on the \emph{whole} genuine cone is
\[
 E(u)=\inf_{f\in C(K),\ f\geq u}\mathscr E(f),\quad
 \E^1_\theta=\{u\in\PSH(\theta):E(u)>-\infty\},\quad B(u)=-E(u)
\]
for normalized $u$. It is concave, monotone, weakly usc, decreasing-net
continuous, has $E(R)=0$ and $E(u+c)=E(u)+cV$, and satisfies
$E(P_\theta(f))=\mathscr E(f)$ for continuous obstacles. Every normalized
$u\in\E^1_\theta$ has decreasing nonpositive rational model majorants
$f_i\downarrow u$ with $f_i(v_0)=0$, whose envelopes $a_i=P_\theta(f_i)$
decrease all-point to $u$ and whose energies decrease to $E(u)$.
For fixed $h\in C(K)$, their perturbed envelopes decrease to
$P_\theta(u+th)$ at each fixed $t$. Outer energy on usc obstacles equals
the energy of this envelope. Its continuous-obstacle value is the
full-section trace constructed from $W_m(f;s)$ below.
\item[S3.] There is an actual Kodaira identity $L=aA+E_0$ with $a>0$,
$A$ ample and $E_0\geq0$. Write $Z=E_0^\an$, $U=K\setminus Z$.
There is a normalized actual $w_Z\in\E^1_\theta$ with $w_Z=-\infty$
on $Z$. Genuine relative contraction holds with a dimension constant
$C_0\geq1$:
\[
 B((1-r)R+ru)\leq C_0r^2B(u),\qquad 0<r<1.
\]
\item[M1.] A target-independent map $M$ on every actual $\E^1_\theta$
has positive original Radon outputs of raw mass $V$, constant invariance,
and $M(P_\theta(f_D))=\mu_D$ from \eqref{finite-mu}.
For the canonical majorants in S2, $M(a_i)\to M(u)$ weakly on $K$.
Every such output gives $Z$ zero mass. The map is not defined using an
inverse target construction.
\item[M2.] For any fixed finite list of normalized nonpositive rational
model obstacles, resolve the \emph{full} $|mL|$ system and their full
relative lattices on common smooth equivariant models. One has the
actual generic identity $p_m^*L=A_m+T_m$, with $mA_m$ generated,
$mT_m$ effective, and bounded relatively nef model scalars $\psi_{m,D}$
on this \emph{same actual} $A_m$, such that
\[
 \phi_{m,D}=q_{T_m}+\psi_{m,D},\quad\phi_{m,0}=q_{T_m},\quad
 A_m^n\to V,\quad e_{A_m}(\psi_{m,D})\to E(P_\theta(f_D)),\quad
 \psi_{m,D}(v_0)\to0.
\]
These full complete-system stages are minorants of their original
envelopes. Their pushed auxiliary raw top measures converge to $\mu_D$.
All degree limits are cofinal divisibility limits with nonempty systems.
The auxiliary bounded nef trace agrees with the finite intersection
energy $e_{A_m}$. Multiplication by $s_{\ell T_m}$ injects
$H^0(X_m,\ell A_m)$ into the \emph{full original} $H^0(X,\ell L)$
in cleared degrees. No assertion of original $\E^1$ membership is made
for the physical functions $q_{T_m}+\psi_{m,D}$.
\end{enumerate}
Here M2 is only a finite-list, finite-obstacle compatibility assertion.
It does not assume an unbounded mixed product, arbitrary approximation
independence, attainment, comparison, whole supporting planes, or separation.
Theorem~\ref{taylor} supplies its finite variation ingredient; no further
uniform Taylor hypothesis is hidden in S--M.

The supplied scalar constructions and moving-product arguments suggest
S--M, but checking all their antecedents and unbounded passages is separate
work. In particular G1 alone does not grant S1, and G0 does not grant M1
on a big nonnef class. The conditional theorem below will show exactly
what follows if this package is established.

\section{Independent bounded-obstacle density and directed integration}
The density result in this section requires only the single-net definition,
genuine obstacle envelopes from S1, and ordinary Kodaira/ample background.
It reconstructs the bounded ingredient of the supplied ROOT198 bridge.
Let $W=q_{E_0}\leq0$. It is finite continuous on $U$, and tensoring its
single section with a weight-zero generating system of $aA$ realizes it
as an actual FS scalar on $L$.

For $f\in C(K)$ and nonzero $s\in H^0(X,mL)$ put
\[
 W_m(f;s)=\inf_K(mf-\log\|s\|),
 \qquad
 \phi_{m,f}=m^{-1}\max_j(\log\|s_{mj}\|+\lfloor W_m(f;s_{mj})\rfloor).
\]
At a section zero the bracket is $+\infty$. Its bounds are
$m\min f\leq W_m(f;s)\leq m f(v_0)$. Superlevel subspaces of the
full section space form a finite closed-endpoint flag; choose an adapted
basis. The ultrametric inequality over trivial $k$ makes $\phi_{m,f}$
basis independent. Powers of sections and superadditivity of weights
make these stages increase in divisibility order. In the model case
they are exactly the complete relative-lattice stages by the all-point
section-extension calculation in G1. Real FS weights are dominated in
the supremum: raising a degree-$r$ section to its $\ell$th power loses
less than $1/(r\ell)$ upon flooring its weight.

\begin{lemma}[Actual bounded density]\label{density}
If $u_f=P_\theta(f)$ is genuine, then
\[
 u_f=\sup_m\phi_{m,f}\quad\text{on }U.
\]
This envelope is finite continuous on $U$ and the stages converge
uniformly on each compact subset of $U$. For $|f|\leq C$, the relative
function $u_f-R$ is bounded by $C$ on $U$.
\end{lemma}
\begin{proof}
Take $C$ a positive integer with $|f|\leq C$. The genuine minorant
$W-C$ makes $u_f$ finite on $U$. Given a defining net $h_i\downarrow u_f$
on actual $B_i$, fix rational $0<t<1$ and $\epsilon>0$ \emph{before}
taking a tail. The actual difference
\[
 L-(1-t)B_i-t(L-aA)
\]
is eventually ample, since its numerical class tends to $ta[A]$.
Compactness of the nested closed sets $\{h_i\geq f+\epsilon\}$ gives
a tail with $h_i<f+\epsilon$ on \emph{all} of $K$, including supported
poles. Tensor with a generated system of that actual ample difference.
The scalar $(1-t)(h_i-\epsilon)+t(W-C)$ is then an actual FS function
on $L$ below $f$ everywhere. At $x\in U$, let $i$ tend to its endpoint,
then $t\downarrow0$, then $\epsilon\downarrow0$. Full-system domination
including the real-weight flooring observation proves the reverse
inequality $u_f\leq\sup_m\phi_{m,f}$. The other inequality follows
because each stage is a genuine minorant.

One cleared Kodaira-section tensor degree has no common zeros on $U$.
For all its multiples the stages are finite continuous there; their
supremum is lsc, while the genuine envelope is usc. Thus equality makes
it continuous. On a fixed compact subset the decreasing continuous gaps
to this envelope have empty positive-margin superlevel intersections.
Compactness yields uniform convergence for the divisibility-directed
stages. Finally constant shifts and envelope order give
$R-C\leq u_f\leq R+C$ on $U$. No equality on $Z$ is asserted.
\end{proof}

\begin{lemma}[Bounded weighted passage]\label{bounded-passage}
If positive Radon measures $\sigma_i\to\sigma$ weakly on $K$ with
converging masses, $\sigma(Z)=0$, and $b$ is bounded continuous on $U$,
then its bounded extension by zero on $Z$ has converging integrals.
The same holds for bounded changing section quotients which converge
uniformly to $b$ on every compact subset of $U$.
\end{lemma}
\begin{proof}
Choose a compact $F\subset U$ with arbitrarily small omitted $\sigma$
mass. By compact Hausdorff normality choose a continuous $0\leq\chi\leq1$,
equal to one on $F$ with compact support in $U$. The product $\chi b$
extends continuously by zero on $K$. Pass that fixed continuous integral
through weak convergence. Bound each omitted part by the common bound
times $\sigma_i(K)-\int\chi d\sigma_i$, which tends to its arbitrarily
small limiting value. For changing quotients first fix this collar,
then their uniform error on its support, and only then take the index.
This proves the claim without indicators of changing Borel sets or
uniform convergence at polar points.
\end{proof}

\begin{lemma}[Radon finite-cover lemma]\label{radon}
Let $c_i\geq0$ be a directed increasing net of continuous functions on
$U$, and let $\nu$ be finite Radon on $K$ with $\nu(Z)=0$.
Then $\int_U\sup_i c_i\,d\nu=\sup_i\int_U c_i\,d\nu$,
allowing $+\infty$.
\end{lemma}
\begin{proof}
For $a>0$, the open set $\{\sup_i c_i>a\}$ is the union of the open
superlevels of the $c_i$. Each compact subset has a finite subcover;
directedness gives one index dominating that finite set. Radon inner
regularity identifies the superlevel mass with the supremum of these
masses. Finitely many threshold masses can be approached simultaneously
by a common index. Use finite lower layer-cake step sums for the cap
$\min(\sup_i c_i,N)$, refine the threshold mesh, then let $N\to\infty$.
The reverse inequality follows from pointwise order. This proves the
net integral assertion; measurable monotone convergence for arbitrary
uncountable nets was not used.
\end{proof}

\section{Reconstruction of arbitrary-target output approximation}\label{target-approx}
For a normalized actual $u\in\E^1_\theta$, all differences below are
taken where $R$ is finite. An intrinsic target is a positive original
Radon measure $\nu$ with
\begin{equation}\label{target}
 \begin{aligned}
 &\nu(K)=V,\qquad\nu\{R=-\infty\}=0,\\
 &I_\nu(u)=\int(R-u)d\nu<\infty\quad\text{for every actual }u\in\E^1_\theta,\\
 &D_\nu=\sup_{u(v_0)=0}\{E(u)+I_\nu(u)\}<\infty.
 \end{aligned}
\end{equation}
The redundant integrability clause is displayed to retain every actual
test: finite dual energy directly forces each integral to be finite.
No absolute integrability of $R$ is imposed. S3 and $w_Z$ give
$\nu(Z)=0$, since its relative cost is infinite on $Z$ outside the
already excluded reference poles. These are intrinsic targets of the
specified reference, not a support-restricted class.

\begin{proposition}[Same fixed increment at every genuine base]\label{fixed-increment}
Under S--M, \eqref{taylor-eq} holds with its same constants at every
normalized genuine $u$ of cost at most $B$, replacing its left side by
\[
 |E(P_\theta(u+th))-E(u)-t\int h\,dM(u)|.
\]
Both derivative signs exist for every original $h\in C(K)$ and equal
$\int h\,dM(u)$.
\end{proposition}
\begin{proof}
Apply Theorem~\ref{taylor} to the S2 model majorants $f_i$. Their
envelope costs are at most $B$. At fixed $t$, the perturbed envelopes
and their energies decrease to those of $P_\theta(u+th)$, and M1
passes the fixed continuous $h$ integral. Thus the whole finite
increment estimate passes with unchanged constants. There is no
exchange of a derivative with a net. Divide by $t$ with both signs,
first let $t\to0$ at fixed $c>B/V$, then $c\to\infty$.
Uniform model approximation of any original continuous $h$ changes
the divided energy increment and the measure integral by at most
$V$ times its sup-norm error. Fix the approximation before $t\to0$,
then let this error tend to zero. This proves both continuous signs.
\end{proof}

\begin{proposition}[Output approximation before attainment]\label{output-approx}
Under S--M, every target \eqref{target} has normalized model-envelope
near-maximizers $u_k$ with
\[
 E(u_k)+I_\nu(u_k)\geq D_\nu-1/k,\quad
 B(u_k)\leq Q_\nu:=4C_0D_\nu+2,\quad M(u_k)\to\nu\text{ weakly}.
\]
\end{proposition}
\begin{proof}
For the actual majorants $a_i$ of any $u$, the costs $R-a_i$ are
nonnegative bounded continuous on $U$ by Lemma~\ref{density}, increase
to $R-u$, and have finite $\nu$ integral. Lemma~\ref{radon} and S2
show convergence of the objective. Its supremum over model envelopes
is therefore exactly $D_\nu$. Choose a sequence within $1/k$ of that
real supremum; this does not impose metrizability on $K$.

Reference contraction gives
$rI_\nu(u)-C_0r^2B(u)\leq D_\nu$. At $r=1/(2C_0)$ it gives
$I_\nu(u)\leq B(u)/2+2C_0D_\nu$. The near-maximizer lower bound
then gives $B(u_k)\leq Q_\nu$.

For a fixed rational model $h$, perturb the \emph{actual base}
$u_k+th$, not its larger physical model obstacle. Its envelope $z_{k,t}$
lies between $u_k\pm|t|\|h\|_\infty$ and is actual finite energy by S2.
The obstacle inequality gives
$I_\nu(z_{k,t})\geq I_\nu(u_k)-t\int h\,d\nu$.
Normalize it only when using \eqref{target}; the raw constant law and
mass keep the objective unchanged. Near maximality implies
\[
 E(z_{k,t})-E(u_k)-t\int h\,d\nu\leq1/k.
\]
Proposition~\ref{fixed-increment} at bound $Q_\nu$, with both signs,
gives for $0<t<t_c$
\[
 \left|\int h\,d(M(u_k)-\nu)\right|
 \leq \frac1{kt}+K_ct/2+
           \|h\|_\infty(2Q_\nu+2V\|h\|_\infty)/c.
\]
Send $k\to\infty$ at fixed $c,t$, then $t\downarrow0$ at that $c$,
then $c\to\infty$. Uniform model density and the common raw mass
extend convergence to every continuous test. This independently
reconstructs the conditional ROOT199 bridge, before attainment.
\end{proof}

\section{Midpoint transfer with the lost-volume error retained}\label{transfer}
For normalized genuine $a,b$ define the scalar midpoint defect
\[
 \delta(a,b)=E((a+b)/2)-(E(a)+E(b))/2\geq0.
\]
The central new original-carrier assertion of this paper is the following
\emph{conditional} absolute estimate.

\begin{theorem}[Absolute midpoint modulus]\label{absolute-original}
Assume S--M. For any normalized $a,b,w\in\E^1_\theta$,
\begin{equation}\label{modulus-output}
 \int_U|a-b|\,dM(w)
 \leq C_n\delta(a,b)^\alpha B(w)^{1/2}
           \max\{B(a),B(b),B(w)\}^{1/2-\alpha},\quad \alpha=2^{-n}.
\end{equation}
Every genuine relative test is integrable for $M(w)$, and these outputs
have finite intrinsic dual energy. For any target \eqref{target},
\begin{equation}\label{modulus-target}
 \int_U|a-b|\,d\nu
 \leq C_n\delta(a,b)^\alpha Q_\nu^{1/2}
          \max\{B(a),B(b),Q_\nu\}^{1/2-\alpha}.
\end{equation}
All these integrals use original actual functions and relative values.
For unbounded pairs the absolute difference is defined where both relative
costs are finite, which has full measure by the integrability proof below
(or by the target condition), and assigned zero on its null complement.
No subtraction of two infinite values is made.
\end{theorem}
\begin{proof}
First fix normalized model envelopes $a=P(f_H)$, $b=P(f_J)$,
$w=P(f_D)$. Use M2's common moving $A_m,T_m$ and set
$\xi_m=\psi_{m,H}$, $\eta_m=\psi_{m,J}$, $\psi_m=\psi_{m,D}$.
The finite endpoint energies and generic values have the stipulated
limits. We need only an \emph{upper} bound for the midpoint auxiliary
energy; its equality to the original midpoint energy is not presumed.

Fix one finite original continuous $f\geq(a+b)/2$. The literal section
minorants give on the entire pulled carrier
\[
 q_{T_m}+(\xi_m+\eta_m)/2\leq(a+b)/2\leq f.
\]
For this fixed $m$, inject the full auxiliary $H^0(\ell A_m)$ via
$s_{\ell T_m}$ into the full original $H^0(\ell L)$.
Every injected section has original $f$-weight at least its auxiliary
midpoint-metric weight: cancel the actual section factor where finite,
and note that its zero branches have original weight expression
$+\infty$ and cannot lower the infimum. Analytic surjectivity identifies
the original infimum with the infimum on all lifts. All remaining
original weights are at least $\ell\min_Kf$, since canonical section
norms are at most one. Extend an adapted subspace basis to a full
adapted basis; the quotient weights have this lower bound.

Take $\ell\to\infty$ with $m,f$ fixed. Full section asymptotics and the
finite auxiliary nef trace from M2 give exactly
\begin{equation}\label{trace-error}
 \mathscr E(f)\geq
 e_{A_m}((\xi_m+\eta_m)/2)+(\min_Kf)(V-A_m^n).
\end{equation}
The rounding error is bounded by the ranks and disappears at this
first degree limit. The last term is generally negative and has not
been discarded at finite $m$. Now let $m\to\infty$; only for this
\emph{fixed finite} $f$ does the term vanish. Finally take the outer
infimum over $f$. Thus
\[
 \limsup_m\delta_{A_m}(\xi_m,\eta_m)\leq\delta(a,b).
\]
This uses no original $\E^1$ membership of a physical complete-system
function and no equality of its energy with an auxiliary energy.

Theorem~\ref{absolute-finite} now bounds the integrals of
$|\xi_m-\eta_m|$ against the auxiliary raw output of $\psi_m$.
Its generic error $|\xi_m(v_0)-\eta_m(v_0)|A_m^n$ tends to zero
\emph{after} this finite estimate; all $J$ factors tend to $B(a),B(b),B(w)$.
At the pushed function-field divisorial atoms, that test is exactly
$|\phi_{m,H}-\phi_{m,J}|$. On $U$ these quotients converge uniformly
on compact sets to $|a-b|$, by Lemma~\ref{density}; their global
support values have a fixed bound from time-lattice comparison.
Indeed the full weights for $f_H$ and $f_J$ differ by at most
$m\|f_H-f_J\|_\infty$, up to a rounding error of one.
Lemma~\ref{bounded-passage}, with limiting finite divisorial $\mu_D$
which misses $Z$, identifies the full limiting integral. This proves
\eqref{modulus-output} for model pairs at model outputs.

For a fixed model pair $a,b$, its absolute difference is bounded
continuous on $U$. Apply Lemma~\ref{bounded-passage} along canonical
model outputs of an arbitrary $w$, using M1 and their bounded costs,
to obtain \eqref{modulus-output} for this pair at $M(w)$.
Apply it instead along Proposition~\ref{output-approx} to obtain
\eqref{modulus-target} for each fixed model pair. This requires only
bounded relative integration off the actual divisor, not optional
unbounded approximants or a changing test against $\nu$.

We next justify all unbounded tests. Take $a=R$ and a decreasing
normalized model-majorant envelope $b_j$ of $v\in\E^1_\theta$.
Since $v\leq b_j\leq R$, monotonicity gives
$\delta(R,b_j)\leq B(b_j)/2\leq B(v)/2$.
The model-pair estimate uniformly bounds $\int(R-b_j)dM(w)$ by
\[
 C_n B(v)^\alpha B(w)^{1/2}\max\{B(v),B(w)\}^{1/2-\alpha}.
\]
These nonnegative continuous costs on $U$ increase to $R-v$.
Lemma~\ref{radon} therefore proves its actual integrability at the
fixed output. In particular the model majorants converge in its
relative $L^1$. The same directed integral argument works at a
target by its prescribed integrability.
For two arbitrary tests $a,b$, pass their two model-majorant nets
separately in $L^1$ at the fixed measure; their midpoint energies
decrease by S2. The absolute-value error is bounded by the sum of
those two $L^1$ errors. This proves both full estimates.

The reference-pole set is contained in $Z$, since $R\geq W$.
Outputs already give $Z$ zero mass. To check finite intrinsic dual
energy, use the displayed moment bound. If $B(v)\geq B(w)$ it is at
most $C_n\sqrt{B(v)B(w)}$; if $B(v)<B(w)$ it is at most $C_nB(w)$.
Subtracting $B(v)$ gives a finite supremum, of size at most a
dimension constant times $B(w)$. This verifies the target condition
for outputs without an inverse definition.
\end{proof}

The argument also reconstructs the optional fixed-test part of ROOT199,
but its hypotheses should be kept explicit. If $b=R-v$ on $U$, extended
by zero on $Z$, then $b$ is nonnegative lsc on $K$.
If actual $a_\eta\in C(K)$ satisfy $0\leq a_\eta\leq b$ and
$\sup_k\int(b-a_\eta)dM(u_k)<\eta$, the nonnegative lsc Portmanteau
inequality gives $\int(b-a_\eta)d\nu\leq\eta$.
Fixed continuous integrals then pass, giving fixed-test weighted convergence.
On compact Hausdorff $K$, this Portmanteau inequality follows by approximating
nonnegative lsc functions from below by continuous functions, using Radon
inner regularity of their open superlevels. It holds for nets.
The existence of these below-approximants is an additional assumption in
that optional argument; the observation alone does not prove it.
Our midpoint proof below does not assume those approximants.

\section{Actual conditional attainment, excluding the energy jump}
We first establish precisely the finite supporting planes used in the proof.

\begin{lemma}[Finite whole planes]\label{finite-plane}
Under S--M, if $a=P_\theta(f_D)$ is a model envelope, then for every
actual $v\in\E^1_\theta$
\begin{equation}\label{plane}
 E(v)\leq E(a)+\int(v-a)d\mu_D.
\end{equation}
This is a finite divisorial integral and is finite for every actual test.
\end{lemma}
\begin{proof}
Concavity and the finite derivative in \eqref{finite-mu} give a plane
at $f_D$ for every continuous obstacle. Uniform model approximation
extends this continuous-obstacle plane to $C(K)$.
Take decreasing continuous majorants of $v$ and pass their values
on the fixed finite divisorial support; G0's actual ample embedding
ensures all of these limiting divisorial values are finite.
S2 passes the outer energies. Thus
$E(v)\leq\mathscr E(f_D)+\int(v-f_D)d\mu_D$.
At $v=a$, S2 gives equality of the two energies, while $a\leq f_D$.
Hence the nonnegative contact integral $\int(f_D-a)d\mu_D$ is zero.
Replacing $f_D$ by $a$ in the plane proves \eqref{plane}.
No passage at nonalgebraic poles occurs at this finite support.
\end{proof}

\begin{theorem}[Conditional attained equation for the outer map]\label{attainment}
Assume S--M for the entire retained big domain. Then every target
\eqref{target} has an actual normalized maximizing $u\in\E^1_\theta$
on $K$, and $M(u)=\nu$. In particular no abstract completion or lift
on a modification replaces the solution. This theorem is conditional
on S--M and does not identify $M$ with every required mixed operator.
\end{theorem}
\begin{proof}
Use Proposition~\ref{output-approx} and put $\epsilon_k=1/k$,
$F_\nu(v)=E(v)+I_\nu(v)$.
Linearity of actual target relative integrals on convex combinations gives
\begin{equation}\label{near-mid}
 0\leq\delta(u_i,u_j)
 =F_\nu((u_i+u_j)/2)-(F_\nu(u_i)+F_\nu(u_j))/2
 \leq(\epsilon_i+\epsilon_j)/2.
\end{equation}
For each rational divisorial $z$, the raw target $V\delta_z$ is
intrinsic: it misses reference poles, every genuine test is finite
there by \eqref{coordinate-bound}, and its normalized relative integral
is at most $VC_z$, so its dual energy is finite. This conclusion
uses no Dirac inverse equation. Applying \eqref{modulus-target}
to this target shows that $u_k(z)$ is Cauchy.
S1 compactness and all-point divisorial determination therefore give
one actual normalized weak limit $u$ for the whole sequence. S2
upper semicontinuity and the cost bound give $E(u)\geq-Q_\nu$,
so the limit is genuinely in $\E^1_\theta$.

The finite planes and \eqref{modulus-output} give
\[
 |E(u_i)-E(u_j)|\leq
 \max_{l=i,j}\int|u_i-u_j|dM(u_l)
 \leq C_{n,Q_\nu}(\epsilon_i+\epsilon_j)^\alpha.
\]
Hence these energies have a real limit $e$ and $E(u)\geq e$.
An energy jump is excluded by a \emph{fixed finite} plane.
For fixed $i$, $M(u_i)=\mu_i$ has finite divisorial support. Weak
convergence of the potentials on that support gives
\[
 \int|u-u_i|d\mu_i
   =\lim_j\int|u_j-u_i|d\mu_i
   \leq C_{n,Q_\nu}\epsilon_i^\alpha.
\]
Lemma~\ref{finite-plane} now gives
$E(u)\leq E(u_i)+C_{n,Q_\nu}\epsilon_i^\alpha$.
Let $i\to\infty$. Thus $E(u)=e$. The direction of this inequality
is the missing direction that compactness and energy usc alone cannot give.

The midpoints $(u_i+u)/2$ converge weakly to $u$. Concavity gives
their energy lower bound tending to $E(u)$; energy usc gives the matching
upper bound. Therefore $\delta(u_i,u)\to0$.
The full target modulus \eqref{modulus-target} yields
$\int|u_i-u|d\nu\to0$. The varying relative integrals, previously
uncontrolled, now converge to the actual limit's integral. Together
with the proved energy equality this gives $F_\nu(u)=D_\nu$.

For fixed original continuous $h$, set $z_t=P_\theta(u+th)$.
The same actual obstacle inequality used earlier gives
$E(z_t)-E(u)-t\int h d\nu\leq0$ by maximality.
Both continuous derivative signs from Proposition~\ref{fixed-increment}
give $\int h dM(u)=\int h d\nu$. Common finite mass and Riesz
uniqueness identify these measures. Attainment precedes this equation;
the target did not define $M$.
\end{proof}

\section{Conditional whole planes and every-solution separation}
The same absolute estimate also addresses an obligation independent of
attainment: every operator solution must belong to the maximizing family,
and that family must be separated by actual probes.

\begin{proposition}[Fixed tests and whole planes]\label{whole-planes}
Under S--M, if $\sigma_i=M(w_i)\to M(w)$ weakly with uniformly bounded
costs, then every \emph{fixed} actual relative finite-energy test has
converging integrals. Moreover
\begin{equation}\label{whole-plane}
 E(v)\leq E(u)+\int(v-u)dM(u),\qquad u,v\in\E^1_\theta.
\end{equation}
\end{proposition}
\begin{proof}
For a fixed normalized $v$, take model majorants $v_j\downarrow v$.
Their ordered midpoint defects satisfy
$\delta(v_j,v)\leq(E(v_j)-E(v))/2\to0$.
Estimate \eqref{modulus-output} therefore gives uniform absolute
$L^1$ error over the bounded-cost family of outputs for replacing
$R-v$ by the bounded relative test $R-v_j$.
At each fixed $j$, Lemma~\ref{density} and
Lemma~\ref{bounded-passage} pass that bounded relative integral through
weak convergence, since the limiting output misses $Z$. This proves
the fixed-test assertion, also after multiplying by a fixed $h\in C(K)$.
The collar is fixed at this bounded step before the output index.

Take canonical model envelopes $a_i\downarrow u$. Their costs are
bounded and their outputs converge by M1. Apply their finite planes
to $v$, writing $v-a_i=(v-u)-(a_i-u)$ on $U$.
The first integral converges by the fixed-test result just proved,
applied separately to $v$ and $u$. The second tends to zero by
\eqref{modulus-output} and
$\delta(a_i,u)\leq(E(a_i)-E(u))/2$.
Pass the decreasing energies to obtain \eqref{whole-plane}.
Constants are restored using raw mass $V$ and the energy constant law.
\end{proof}

\begin{corollary}[Every-solution uniqueness for the conditional map]
Under S--M every actual $M$-solution for an intrinsic target is unique
modulo a real constant. Every such solution maximizes the intrinsic objective.
\end{corollary}
\begin{proof}
If $M(a)=\nu$, \eqref{whole-plane} gives $F_\nu(v)\leq F_\nu(a)$
for every actual $v$. Normalize any two solutions $a,b$ at $v_0$.
They maximize the same objective; its value at their midpoint is
at least this common maximum by concavity and at most it by definition.
Thus $\delta(a,b)=0$.
For every rational divisorial $z$, apply the intrinsic-probe estimate
\eqref{modulus-target} at $V\delta_z$. It gives $a(z)=b(z)$.
G0 order determination in the common actual ample cone gives equality
on \emph{all} original points, including supported poles. Removing
normalization restores one real constant. This proof treats every
solution of the conditional map; it does not merely pick one maximizer.
\end{proof}

This separation mechanism does not supply strict Borel locality or comparison.
No claim that $\{u<v\}$ is open appears in it. Nor does it identify every
unbounded mixed moving quotient with an original genuine test. Thus even
the conditional attained-and-separated equation above is not the full
original P16.5 theorem required by the assignment.

\section{The additional zero-volume range actually proved}
Here let the actual rational bundle $L$ be numerically trivial on the
same smooth integral projective $X$, including nontorsion elements of
$\operatorname{Pic}^0(X)$. Use the literal varying-actual-bundle cone
at numerical class zero and the canonical zero scalar as reference.
The following result uses G0 and the single-net definition directly,
and does not require S--M.

\begin{proposition}[Numerical zero]
In this convention $\PSH(0)=\mathbb R$. Its raw energy is zero,
$\E^1(0)=\mathbb R$, and its raw top operator on constants is zero.
The only positive mass-zero target is literally the zero measure,
and its equation is attained and unique modulo constants.
\end{proposition}
\begin{proof}
For every positive rational $\epsilon$, the actual $L+\epsilon A$
is ample. Its weight-zero generating system is zero on every original
point. These constant scalars give one defining net for zero, and
constant shifts give all real constants. Conversely normalize
$u\in\PSH(0)$ to generic value zero. Integer tensoring makes
$m u\in\PSH(0)$ for every positive integer $m$. Embed each into
one fixed actual ample cone by the elementary embedding in Section 1.
The same ample coordinate bound gives $-C_z\leq m u(z)\leq0$ for
every $m$, hence $u(z)=0$ for each rational divisorial $z$.
G0 all-point determination gives $u=0$ everywhere.

The constant raw ample energies are $c\epsilon^n(A^n)$ and the
constant raw top measures have this positive mass. Positivity bounds
their integral against any continuous test by its sup norm times
$\epsilon^n(A^n)\to0$. Their unique raw limit is zero, and the
energy is zero on every constant. This gives the stated energy and
operator on the \emph{entire} classified cone. A positive Radon
measure of zero total mass is zero. Its relative dual energy is zero.
The quotient of solutions by constants has one element. Comparison
and both zero energy derivative signs are immediate. No zero volume
is divided out and no probability measure is retained.
\end{proof}

This is precisely the numerical-zero range, not the full pseudoeffective
zero-volume cone. On $\mathbb P^1\times\mathbb P^1$ take the nonzero nef
$L=p_1^*\OO(1)$. Its volume is zero, and the pullback of
$\max(\log\|z_1\|_{\rm triv},-1)$ is a nonconstant bounded generated-system
scalar in its nef cone, while zero is another. The same scalar can be
written with the actual generated sections $z_1,z_0$ and respective
weights $0,-1$, so it is not a merely local formula.
Tensor into $L+\epsilon\OO(1,1)$. Their positive raw top masses are
$2\epsilon(1+\epsilon)\to0$, hence both raw top limits vanish although
the original potentials differ by more than a constant. Thus a general
raw mass-zero top equation cannot retain uniqueness on that entire
nonzero nef cone. This is a safeguard, not a counterexample to the
requested positive-volume theorem or a solution for all zero-volume data.

\section{Precise unproved statements and canonical result}\label{gaps}
The original completion criterion is not achieved. The retained domain
remains \emph{every} actual big rational bundle on every retained smooth
integral projective $n$-fold over every retained field. The precise boundary
of this submission is as follows.
\begin{enumerate}
\item Independently establish S1--S3 on the literal single-net cone,
including closure across old supported poles, original envelopes, actual
normalized compactness, weak usc and all-point outer full-section energy,
canonical nonpositive model majorants, finite-energy divisor poles, and
genuine contraction. None is a consequence just of deleting ``ample''
from G0; G1's one endpoint does not grant the whole package.
\item Independently establish M1--M2 in the entire domain. In particular,
complete the target-independent output construction, its finite-model
compatibility, every required weak limit, algebraic-divisor nullness,
the actual shared horizontal-factor identities and the finite-moving
energy and raw top-measure limits. The finite Taylor calculation in
Section~\ref{variation-section} is not this whole construction.
\item Construct and identify the \emph{genuine original} finite-energy
operator with every admissible approximation, all unbounded mixed
products, all raw weighted limits and their independence. The finite
absolute modulus and the conditional whole planes do not identify
simultaneously changing physical quotients which are outside original
$\E^1$. Ordinary big positive products are not made multilinear here.
\item Prove original strict-winning locality and strict Borel comparison
for every actual finite-energy pair. The set $\{u<v\}$ need not be
open. Neither the fixed-continuous weak limits nor the discarded one
algebraic divisor justify passing its indicator through changing measures.
Nonalgebraic negligible-set nullness needs its own proof when used.
\item To make the conditional attained-and-separated equation in
Theorem~\ref{attainment} an original-scope theorem, first discharge
the package in the preceding items and identify the map used there
with the genuine required operator. The attainment and strict-probe
arguments themselves have full proofs here under their explicit
antecedents; they are not independent grants for those antecedents.
\item Specify and prove any further nonzero zero-volume pseudoeffective
range with its actual reference and replacing datum. Only the
numerically trivial rational-bundle range is established here.
\end{enumerate}
These are unresolved assertions of this draft, not a bibliographical
claim that no theorem exists. No other P16 family axis is claimed solved.
The exact unchanged original task and full contextual family are preserved
in the frozen input. No narrowed endpoint or formal conditional transfer
is counted as the original solution. The canonical result is
\texttt{candidate=false}.

\appendix
\section{Primary evidence, literal reading coordinates, and authority}
The two caller grants are preserved whole in
\texttt{supplied/unit-00.txt} and \texttt{supplied/unit-01.txt}.
Every other supplied body and ledger is preserved as mathematical data
in \texttt{supplied/}; \texttt{supplied/index.json} maps each whole unit
to the inclusive LF coordinates of the exact decoded assignment.
No historical path from those ledgers was opened. No independent reviewer,
private output, other campaign, or earlier claimed acceptance was inspected.

The raw public primary connector responses are saved in \texttt{sources/}.
The script \texttt{extract\_sources.py} constructs six complete pertinent
statement/proof units and maps their local LF lines to literal renderer
labels, preserving PDF page labels separately. Its corrected parser handles
multiple source documents inside one connector text block. All six selected
intervals are complete and overlap consistently; no missing renderer label
or conflicting text remains. Local coordinates mean UTF-8 decoding followed
by splitting on LF, inclusive and one-based. Renderer labels are zero-based;
they are not represented as LF lines of a locally downloaded original PDF.
No original PDF download or original PDF byte hash is asserted.

\begin{longtable}{p{.16\linewidth}p{.24\linewidth}p{.51\linewidth}}
Version & Exact retained coordinates & Application and limitation\\\hline
BJ v4, arXiv:\allowbreak 1801.08229v4, PDF date 24 March 2022 &
Renderer 2705--2789, local LF1--85 &
Complete finite Hodge/concavity unit. Finite model calculus only;
the algebraic Hodge input cited there is ordinary background.\\
BJ v4 & Renderer 2790--3107, local LF1--318 &
Complete finite distance, slot replacement and quasi-triangle proof unit.
We reconstruct it with the independently proved $n-j$ orientation.
The printed repeated-potential quasi-triangle expression is not used.\\
BJ v4 & Renderer 6157--6257, local LF1--101 &
Complete ample four-point lemma/proof, diagnostic only. Its unbounded
ample regularization is not transferred to a big nonnef class.\\
BFJ v2, arXiv:\allowbreak math/\allowbreak 0608260v2, posted 4 September 2007 &
Renderer 310--441, local LF1--132 &
Complete positive-product construction, directed compactness, monotonicity,
continuity and nef agreement units. Used in actual smooth big algebraic
determinations, paired through fixed nef tests.\\
BFJ v2 & Renderer 445--545, local LF1--101 &
Concavity/volume statements and complete pertinent volume derivative proof.
Both derivative signs and continuity at the big boundary are retained;
Fujita, section asymptotics and the cited Morse background are explicitly
ordinary algebraic inputs, not an NA equation.\\
Addendum v2, arXiv:\allowbreak 2206.07183v2, PDF date 7 September 2023 &
Renderer 128--257, local LF1--130 &
Complete bounded FS-envelope and model/lattice statements and proofs.
Source matching only for the independently reconstructed density; the
single-net proof does not import its general nonample closure call.
Actual section-extension use is covered by G1's exact signature.\\
\end{longtable}

The BFJ HTML's automatically generated date is not a new source version.
Its characteristic-zero algebraically closed field statement is retained
from the independent primary read, separately from its theorem units.
The source-application ledger distinguishes G0/G1, ordinary published
algebraic premises, new proofs, explicit conditional inputs, and diagnostic
readings. A source ledger from an earlier author is not our reading evidence.
No snippets or unsupplied mathematical verdicts are used as premises.

\section{Reproducibility and checks}
\texttt{frozen/input-original.json} preserves the exact supplied input bytes;
the decoded task and assignment are separately frozen. The complete assignment
still contains every supplied positive-volume and additional zero-volume
statement, without relying on this paper's paraphrases. Run
\texttt{python3 rerun.py} here. It verifies frozen preservation, every
supplied whole-unit coordinate and hash, source-unit maps, source overlap,
backwards local note support, and finite exact rational transcription checks.
It then compiles \texttt{paper.tex} twice using the existing pdfLaTeX.
The program retains argv, stdout/stderr, status, versions, finite check inputs
and outputs, and an artifact manifest. No package is installed.

The finite checks cover dimensions 1--12 for the corrected polynomial and
midpoint coefficients, and concrete rational weighted graphs for the absolute
contraction. These are \emph{finite checks}, not an exhaustive theorem over
all models. The all-dimension arguments are the mathematical proofs in
Sections 2--3. Compiling this paper and passing these tests certify neither
the algebraic imported background nor the conditional package.

No subagents or campaigns were launched. No credentials, account, auth,
runtime configuration, or environment variables were read or changed.
Only this fresh workspace and public primary sources were used.
No operator-provided supervised compute facility was exposed; no
resource-heavy local calculation was launched. The exact native model
identifier is not exposed in the recorded command environment, and no
model/runtime change was attempted. This is an unverified mathematical draft.

\begin{thebibliography}{9}
\bibitem{BJ} S. Boucksom and M. Jonsson,
\emph{Global pluripotential theory over a trivially valued field},
\href{https://arxiv.org/pdf/1801.08229v4}{arXiv:1801.08229v4}.
\bibitem{BFJ} S. Boucksom, C. Favre and M. Jonsson,
\emph{Differentiability of volumes of divisors and a problem of Teissier},
\href{https://arxiv.org/html/math/0608260v2}{arXiv:math/0608260v2}.
\bibitem{ADD} S. Boucksom and M. Jonsson,
\emph{Addendum to global pluripotential theory over a trivially valued field},
\href{https://arxiv.org/pdf/2206.07183v2}{arXiv:2206.07183v2}.
\end{thebibliography}
\end{document}
